%% sections/i-stakeholder.tex - H. R. 9510 Bill v4.0 (full bill) - Appendix I
%% Genre: stakeholder engagement and outreach plan. Source deliverable 09. Table 8.
\billsec{APPENDIX I.}{STAKEHOLDER ENGAGEMENT PLAN (NON-OPERATIVE).}\phantomsection\label{toc:aI}

\uscprov{1}{\textit{Independent research draft. Not an enacted law, not pending
legislation, and not legal advice; not endorsed by the FDA, HHS, any standards body, or
any Member of Congress. Bill: H. R. 9510 (illustrative number; the Clerk assigns the
real number at introduction), 119th Congress, 2d Session. Refer to: House Committee on
Energy and Commerce (Subcommittee on Health). Scope: the pre-introduction engagement
step that precedes any sponsor's introduction.}}

\uscprov{1}{\textbf{Goals.} The operative rule of H. R. 9510 is that no robot-patient
interaction code may be generated or executed in a Physical AI oncology investigation
unless an automated verification, validation, and uncertainty quantification process,
bound to named external consensus standards, has first cleared. Because the gate is
bound to specific published standards, the bodies that author and maintain those
standards are the natural technical reviewers of whether the bindings are stated
correctly. Outreach therefore has four goals: test the standards bindings for accuracy
and currency with the people who wrote them; surface implementation burden and
feasibility concerns early, while the text can still change; build informed, voluntary
support among the communities the bill affects; and keep a faithful, traceable record
of who was consulted and what they said.}

\uscprov{1}{\textbf{Principles.} Five principles govern every contact. First, accuracy
before advocacy: lead with the text and the evidence, not with persuasion. Second, no
claimed endorsement: no group has endorsed this bill, and no message, slide, or email
may imply otherwise. Third, listen to revise: treat objections as inputs to the draft,
not obstacles to it. Fourth, transparency of origin: the draft was AI-assisted, and
that is disclosed up front so a named human owner can take responsibility for the final
text. Fifth, respect for role: standards bodies set consensus standards and do not
legislate, agencies administer the law and do not author bills, and patient groups
speak for patients; the plan never asks any party to act outside its role
\cite{anthropic-claude-code}.}

\uscprov{1}{\textbf{Sequencing.} Engagement runs as a sequenced wave inside calendar
year 2026, because the 119th Congress, 2d Session ends January 3, 2027, and a draft not
introduced and acted on by then must be reintroduced. Preparation comes first: finalize
the briefing package, confirm each standards citation against the most current edition,
and stand up the tracking log. Technical review with the standards bodies follows, with
ASME approached first because ASME V\&V 40-2018 and the in-development VVUQ 70 are the
keystone bindings, then the medical-device and robot-safety bodies, then the remaining
standards organizations; these contacts are corrective in purpose, confirming the
bindings before anyone is asked to support the text. Agency and community briefings come
next, in parallel: FDA (CDRH) and HHS on administrability, and patient-advocacy,
oncology, industry, and academic groups on the substance. Consolidation folds feedback
into a single revised text and reconfirms that the bindings still resolve. Handoff
provides the consolidated text and the engagement record to a prospective House
sponsor's office, with all disclaimers intact. Standards-body review precedes the
support-building briefings on purpose: the bill stands or falls on whether its external
anchors are stated correctly, so that has to be verified before the text is socialized
more broadly \cite{asmevv40, asmevvuq70}.}

\uscprov{1}{\textbf{Risks and respect.} The principal risk is implied endorsement: the
bill names many organizations, and any of them could reasonably worry that being briefed
looks like backing the bill. Mitigation is that every contact opens and closes with the
standing disclaimer, the deck carries it on the title slide, and follow-up email repeats
that no endorsement is sought or implied. A second risk is misstating a standard, which
would be both inaccurate and disrespectful to the body that authored it; mitigation is
the corrective framing above, where standards bodies are asked to correct the bindings
before the text is advocated, plus a final pre-handoff citation check. A third risk is
overstepping role, for example asking an agency to take a position on pending-style
legislation it should not comment on; mitigation is to confine each ask to what fits
that party's role. A fourth risk is opacity about origin; mitigation is to disclose the
AI-assisted drafting and the named human owner in the first contact. A final risk is
fatigue from repeated outreach; mitigation is a single consolidated package, one primary
contact per organization, and no second approach unless invited. Throughout, disagreement
is recorded faithfully and reflected in the text rather than argued away.}

\uscprov{1}{Table 8 is the stakeholder register: for each stakeholder it records the
principal interest, the message carried to that audience, and the engagement phase in
which the contact occurs. The message column is consistent with the not-endorsed-by
framing; the request to a standards body is technical accuracy, not support.}

{\footnotesize
\begin{xltabular}{\textwidth}{@{}L{2.7cm} L{4.0cm} Y L{1.7cm}@{}}
\hline\noalign{\vskip 0.04cm}
\textbf{Stakeholder} & \textbf{Interest} & \textbf{Message} & \textbf{Phase}\\
\hline\noalign{\vskip 0.04cm}
\endfirsthead
\hline\noalign{\vskip 0.04cm}
\textbf{Stakeholder} & \textbf{Interest} & \textbf{Message} & \textbf{Phase}\\
\hline\noalign{\vskip 0.04cm}
\endhead
\hline\noalign{\vskip 0.04cm}
\endlastfoot
FDA (CDRH) and HHS & Administrability of a new \S 515D (21 U.S.C. 360e-5) and its fit
with premarket review and predetermined change control plans (\S 515C) & A sequencing
rule that complements, not duplicates, the existing PCCP and quality-system framework;
the ask is the administrability read a new \S 515D would require & Agency\\
\hline\noalign{\vskip 0.04cm}
Patient-advocacy and oncology groups & Patient safety, informed consent, and equitable
access in trials with embodied AI & A patient-protective rule that requires verification
before an irreversible motion; where should disclosure and protection be stronger &
Community\\
\hline\noalign{\vskip 0.04cm}
Medical-device and AI industry & Feasibility, cost, timelines, and interoperability with
existing pipelines & A predictable, standards-anchored rule with a realistic transition
window; help confirm the transition periods are workable & Community\\
\hline\noalign{\vskip 0.04cm}
Academic and professional societies & Scientific rigor, reporting completeness, and
evidence quality & A rule grounded in the consensus methods your members publish; help
keep the methodology current & Community\\
\hline\noalign{\vskip 0.04cm}
ASME (V\&V 40, VVUQ 70) & The keystone bindings: ASME V\&V 40-2018 credibility
commensurate with risk, and VVUQ 70 machine-learning V\&V in development & The gate
operationalizes the credibility-commensurate-with-risk logic your committees defined;
please confirm the references and VVUQ 70's development status & Technical\\
\hline\noalign{\vskip 0.04cm}
IEC (80601-2-77, 60601-1, 62304) & Medical-device and software safety: essential
performance, single-fault, and software life cycle & The gate protects the essential
performance your standards already define for surgical robots; confirm the editions and
amendments match & Technical\\
\hline\noalign{\vskip 0.04cm}
ISO (13482, ISO/TS 15066, 14971, 13849-1, 9283) & Robot safety and metrics: human-contact
force limits, risk management, control safety, and pose accuracy & The gate verifies the
measurable thresholds your standards specify for robot-patient contact; confirm the
bindings are correctly scoped & Technical\\
\hline\noalign{\vskip 0.04cm}
IEEE (7009-2024) & Fail-safe autonomy and hand-back-to-human escalation & The bill's
escalation logic mirrors your fail-safe design standard; confirm the escalation rule is
consistent with 7009-2024 & Technical\\
\hline\noalign{\vskip 0.04cm}
NASA (NASA-STD-7009A) & Uncertainty quantification as a first-class credibility factor &
UQ is a named, first-class gate criterion, consistent with your models-and-simulations
standard; confirm UQ is treated as your standard intends & Technical\\
\hline\noalign{\vskip 0.04cm}
UL (UL 4600) & Safety case for autonomous products & The gate produces and checks the
structured claim-argument-evidence safety case your standard frames; confirm the
artifact aligns with UL 4600 & Technical\\
\hline\noalign{\vskip 0.04cm}
AMA (CPT Appendix S) & Coding and the assistive, augmentative, autonomous AI taxonomy &
The bill's autonomy vocabulary is consistent with your AI taxonomy; confirm the framing
aligns with Appendix S & Technical\\
\hline
\end{xltabular}}
\vspace{-0.8cm}
\figcaption{Table 8. Stakeholder register for H. R. 9510: the principal interest, the
message carried, and the engagement phase for each audience. Standards-body contacts are
corrective and seek technical accuracy, not endorsement; no group has endorsed this
measure \cite{iso13482, isots15066, iso14971, iec8060127, ieee7009, nasastd7009, ul4600,
cpt-appendix-s}.}
